
Soit alors @@M0@@. Le @@M1@@-homomorphisme @@M2@@ est injectif d'après 3.1. Par conséquent, on peut, pour montrer les relations V 2.4.2, se placer dans @@M3@@. Comme dans @@M4@@ les éléments @@M5@@ et @@M6@@ se décomposent en somme de classes de faisceaux inversibles, il suffit, par additivité, de montrer ces mêmes relations pour des éléments @@M7@@, classes de faisceaux inversibles ; elles peuvent encore s'écrire~:
@@M8@@
en les mettant sous la forme V 2.4.1. Elles ne sont alors autres que les relations V 2.3.1 et 2.3.3 qui définissent la structure de @@M9@@-anneau de @@M10@@.

Indiquons le résultat plus général suivant~:

\textbf{Théorème 3.3.} \uline{Soit} @@M11@@ \uline{un topos commutativement annelé}; \uline{on désigne par} @@M12@@ \uline{le groupe de Grothendieck des} @@M13@@-\uline{Modules localement libres} de rang borné sur @@M14@@. \uline{Alors}, @@M15@@, \uline{muni de sa pré}-@@M16@@-\uline{structure canonique}~(V 2.2 b)), \uline{est} \uline{un} @@M17@@-\uline{anneau}.

Pour montrer que @@M18@@ est un @@M19@@-anneau, il faut montrer les relations V 2.4.1, et par additivité, il suffit de les montrer pour des éléments @@M20@@ qui sont classes de @@M21@@-Modules localement libres de rang borné. Comme les rangs de @@M22@@ et @@M23@@ sont bornés, on peut décomposer @@M24@@ en «~somme disjointe~» finie d'ouverts sur lesquels @@M25@@ et @@M26@@ sont de rang constant. En utilisant la commutation du @@M27@@ aux sommes directes finies, et en mettant les relations V 2.4.1 sous la forme V 2.4.2, on voit qu'il suffit de montrer les relations V 2.4.2, dans chacun des @@M28@@ des ouverts considérés, c'est-à-dire qu'il suffit de montrer les relations V 2.4.2 pour des éléments @@M29@@, classes de @@M30@@-Modules @@M31@@ localement libres de rang constant.

